\section{数据分析与建模}

% 写作提示：先用图表和描述统计展示数据证据，再建立与研究问题匹配的模型；报告估计值、区间或效应，
% 检查模型假设、残差、共线性、过拟合等问题，并解释实际含义。下方模型和数值仅为示例。

\subsection{相关性分析}
计算Pearson相关系数矩阵，分析变量间的线性相关关系。

\subsection{模型构建}
构建多元统计模型，$R^2$为0.896，调整$R^2$为0.887。

\subsection{模型诊断}
进行残差分析，检验正态性、独立性和同方差性假设。
